\begin{proof}[Proof of \cref{obj:lem:selected-contour-perturbation-bound}]
On the selector event, choose the selected bank index \(j\) and the decoded integer
\(N\ge1\).  Write \(C_j\) for the selected circle and set
\[
\widehat T_j
=
\Re\left[
\{N\,2\pi i\}^{-1}
\oint_{C_j}
\frac{\widehat G_{1,n}(z)}{\widehat F_{1,n}(z)}\,dz
\right],
\qquad
T_j
=
\Re\left[
\{N\,2\pi i\}^{-1}
\oint_{C_j}
\frac{G(z)}{F(z)}\,dz
\right].
\]
The exact-identification hypothesis for the selected decoded contour gives
\(\vartheta=T_j\).

For each \(z\in C_j\), the selector event gives
\[
\mu\le |F(z)|,
\qquad
|G(z)|\le C_G,
\qquad
|\widehat F_{1,n}(z)-F(z)|\le d,
\qquad
|\widehat G_{1,n}(z)-G(z)|\le e .
\]
Together with \(d\le\mu-\nu\), the triangle inequality yields the empirical
denominator margin
\[
|\widehat F_{1,n}(z)|
\ge |F(z)|-|\widehat F_{1,n}(z)-F(z)|
\ge \mu-d
\ge \nu .
\]
Since \(\mu>0\) and \(\nu>0\), both denominators are nonzero on \(C_j\).  Hence
\[
\frac{\widehat G_{1,n}}{\widehat F_{1,n}}-
\frac{G}{F}
=
\frac{\widehat G_{1,n}-G}{\widehat F_{1,n}}
+
\frac{G(F-\widehat F_{1,n})}{\widehat F_{1,n}F},
\]
and the two terms have moduli bounded by \(e/\nu\) and
\(C_Gd/(\mu\nu)\), respectively.  Thus, pointwise on \(C_j\),
\[
\left|
\frac{\widehat G_{1,n}(z)}{\widehat F_{1,n}(z)}-
\frac{G(z)}{F(z)}
\right|
\le
\frac{e}{\nu}+\frac{C_Gd}{\mu\nu}.
\]

The contour-integrability hypothesis for the selected index permits integration
of this pointwise bound around the circle of radius \(\rho_j\).  Therefore
\[
\left|
\oint_{C_j}
\frac{\widehat G_{1,n}(z)}{\widehat F_{1,n}(z)}\,dz
-
\oint_{C_j}
\frac{G(z)}{F(z)}\,dz
\right|
\le
2\pi\rho_j
\left(\frac{e}{\nu}+\frac{C_Gd}{\mu\nu}\right).
\]
Multiplication by \(\{N2\pi i\}^{-1}\) and taking real parts give
\[
|\widehat T_j-\vartheta|
\le
\frac{\rho_j}{N}
\left(\frac{e}{\nu}+\frac{C_Gd}{\mu\nu}\right).
\]

It remains to relate this contour value to the returned statistic.  For the same
selected \(j\) and decoded \(N\), the construction of
\(\widehat\theta_{\mathrm{spec},n}\) in
\cref{obj:def:adaptive-contour-estimator} uses the rational midpoint \(y\) of the
real coordinate of the certified evaluation rectangle.  That rectangle contains
the normalized empirical contour value whose real part is \(\widehat T_j\), and
its real-coordinate width is at most \(2/\max\{n,1\}\).  Hence
\[
|y-\widehat T_j|\le \frac{1}{\max\{n,1\}}.
\]
The statistic is the clipping of this rational value to
\([-C_{\theta},C_{\theta}]\).  Since \(|\vartheta|\le C_{\theta}\), projection
onto this interval is contractive relative to \(\vartheta\), so
\[
|\widehat\theta_{\mathrm{spec},n}-\vartheta|
\le |y-\vartheta|
\le |y-\widehat T_j|+|\widehat T_j-\vartheta|.
\]

By the bank construction in \cref{obj:def:contour-bank-handle}, the selected
radius satisfies \(\rho_j<R_1\), and the decoded count satisfies \(N\ge1\).  Using
\(d,e,C_G\ge0\) and \(\mu,\nu>0\),
\[
\frac{\rho_j}{N}
\left(\frac{e}{\nu}+\frac{C_Gd}{\mu\nu}\right)
\le
R_1\max\{\nu^{-1},\,C_G(\mu\nu)^{-1}\}(d+e).
\]
Combining the preceding displays proves the asserted bound.
\end{proof}
